
\section{Summary and Conclusion}\label{sec:security-summary}

The security framework developed for CFI-Care establishes a robust, \textbf{defense-in-depth} architecture that addresses the unique challenges of healthcare information systems. 
By distributing security concerns across the API gateway, a specialized Backend-for-Frontend, and a centralized Identity Provider, the system ensures that \ac{PHI} is protected at every layer of the stack.

The implementation successfully achieves several critical security milestones:
\begin{itemize}
    \item \textbf{Unified Identity Management:} Through the integration of Keycloak, the platform manages complex user lifecycles, including automated \ac{FHIR} resource provisioning via custom event 
    listeners and secure email verification SPIs.
    \item \textbf{Multi-Client Authentication:} The architecture provides seamless security continuity,
     supporting gateway-managed \ac{OIDC} sessions for web browsers and native, hardware-backed token storage for mobile devices via PKCE-enhanced flows.
    \item \textbf{Granular Authorization:} Beyond standard role-based access, the system introduces a dynamic patient-to-doctor consent mechanism, 
    allowing for time-bound, revocable access to sensitive medical records.
    \item \textbf{Operational Integrity:} Hardened infrastructure configurations, including TLS termination,
     security-focused HTTP headers,  a Redis-backed audit trail and centralized secrets management using Vault agents ensure the platform meets modern standards for accountability, non-repudiation and secrets isolation.
\end{itemize}

At its current state, CFI-Care provides a functional and secure foundation ready for clinical validation. 
The identified future work -specifically the integration of standard-based consent documents and independent resource-level enforcement on the \ac{FHIR} server- will serve to further align the platform with international healthcare security regulations and large-scale production requirements.



% CFI-Care implements a defense-in-depth security architecture across multiple
% layers:


% \paragraph{Browser Access (Gateway OIDC + BFF Session Metadata)}
% Browser authentication is enforced at nginx/\texttt{oauth2-proxy}. The BFF
% keeps minimal session metadata for user context and global logout coordination.

% \paragraph{API Layer (Middleware + Validation)}
% Every protected API request flows through bearer JWT verification, then
% scope/role enforcement.
% % Input validation prevents injection attacks. 
% Audit logs capture authentication events for forensic analysis.


% \paragraph{Current State}
% The implementation is functional and ready for incremental hardening:
% \begin{itemize}
%     % \item All authentication endpoints are secured with rate limiting and input validation.
%     \item Browser authentication continuity is maintained through gateway-managed OIDC
%           session handling.
%     \item Mobile authentication is implemented with Authorization Code + PKCE, secure
%           token storage, and refresh/retry handling.
%     \item Registration provisioning is automated via Keycloak Event Listener to BFF
%           provisioning endpoint.
%     \item Per-session and global logout are fully functional.
%     \item Audit trails are captured and retained for 90 days.
% \end{itemize}